% Partie : histoire
% Chapitre : biographies
% Section : Ernst Zermelo

\documentclass[../../../include/open-logic-section]{subfiles}

\begin{document}

\olfileid[fr]{his}{bio}{zer}

\olsection{Ernst Zermelo}

\olphoto{zermelo-ernst}{Ernst Zermelo}

Ernst Zermelo naquit le 27 juillet 1871 à Berlin, en Allemagne. Il
avait cinq sœurs, mais la famille était de santé fragile et seules
trois d'entre elles atteignirent l'âge adulte. Ses parents moururent
eux aussi alors qu'il était jeune : à dix-sept ans, lui et ses sœurs
étaient orphelins. Zermelo s'intéressait vivement aux arts,
en particulier à la poésie. Il était réputé pour sa vivacité d'esprit,
son humour et son esprit critique. Parmi ses réalisations
mathématiques les plus célèbres figurent l'introduction de l'axiome
du choix, en 1904, et son axiomatisation de la théorie des ensembles,
en 1908.

À l'université, Zermelo avait des centres d'intérêt variés et suivit
des cours de physique, de mathématiques et de philosophie. Sous la
direction de Hermann Schwarz, il acheva en 1894, à l'université de
Berlin, sa thèse \emph{Recherches sur le calcul des variations}.
En 1897, il décida de poursuivre ses études à l'université de
G\"{o}ttingen, où les travaux de David
Hilbert sur les fondements des mathématiques l'influencèrent
profondément. En 1899, il obtint son habilitation à Göttingen, où
il enseigna comme privat-docent. Il y fut nommé professeur
titulaire en 1905.

En 1910, Zermelo devint professeur ordinaire à l'université de
Zurich. La tuberculose le contraignit à prendre sa retraite en 1916. Après son rétablissement, l'université de Fribourg
lui confia un poste de professeur honoraire en 1926. Il travaillait
alors sur les fondements des mathématiques. Les travaux de Thoralf
Skolem et de Kurt G\"{o}del l'irritaient, et il critiqua publiquement
leurs méthodes dans ses articles. En 1935, dénoncé pour avoir refusé
le salut hitlérien, il renonça à sa chaire honoraire à Fribourg.
%

Les dernières années de sa vie furent marquées par l'isolement.
Après son départ en 1935, Zermelo s'installa à la campagne, où
il vécut modestement. Il rédigea encore en 1937 une note sur le
paradoxe de Skolem. Réintégré à l'université de Fribourg en 1946,
il fut dispensé d'enseignement. Il se maria en
1944 et devint entièrement dépendant de son épouse à mesure qu'il
perdait la vue. En 1951, il était devenu complètement aveugle.
Il mourut à G\"{u}nterstal, en Allemagne, le 21 mai 1953.

\begin{reading}
Pour une biographie complète de Zermelo, voir \citet{Ebbinghaus2015}.
Ses articles fondamentaux de 1904 et de 1908 peuvent être lus dans
leur version allemande originale
\citep{Zermelo1904,Zermelo1908}. Ses œuvres complètes, y compris ses
écrits sur la physique, sont disponibles en traduction anglaise dans
\citep{Ebbinghaus2010,Ebbinghaus2013}.
\end{reading}

\end{document}
